
\lesson{20/10}
So massless particles do not have $2s+1$ degrees of freedom, but only
the $2$ given by helicity, which in this massless case coincides with chirality.

The value of helicity is \textit{not a Lorentz invariant}, since it depends
on the momentum $\mathbf{p}$.

\subsection{The spin operator}
Combining, in the Weyl basis
$$\gamma_5\sigma^{\mu\nu} = \frac{i}{2}\varepsilon^{\mu\nu\rho\sigma}
\sigma_{\sigma\rho}$$
and, recalling the Pauli-Lubanski vector, it is true that
$$W^\mu = \frac{1}{2}\varepsilon^{\mu\nu\rho\sigma}P_\nu M_{\rho\sigma}=
-\frac{i}{2}\gamma_5\sigma^{\mu\nu}P_\nu.$$
In the rest frame of the particle, the 4-vectors have value
\begin{gather*}
P^\nu = (m,\mathbf{0}),\; W^\mu = (0, -mS^i)
\end{gather*}
and the spin operator is defined to be
$$\Sigma^i = \frac{-2}{m}W^i=\begin{pmatrix}
    \sigma^i&\mathbb{O}\\
    \mathbb{O}&\sigma^i
\end{pmatrix}
$$
With this definition, helicity is
$$h = \frac{\Sigma_iP^i}{|\mathbf{P}|} = \sigma_i\hat P^i$$
and, \textbf{for massless particles only}, chirality is
$$\gamma_5 = \frac{\Sigma_iP^i}{|\mathbf{P}|}$$

\subsection{Solution to the Dirac Equation}
Before solving the equation, it is necessary to introduce the
\textbf{Dirac representation} of the spinor, in which $\gamma^0$ and
$\gamma_5$ have different forms:
$$\gamma^0 = \begin{pmatrix}
    \mathbb{I}&\mathbb{O}\\
    \mathbb{O}&-\mathbb{I}
\end{pmatrix},\; \gamma_5=\begin{pmatrix}
    \mathbb{O}&\mathbb{I}\\
    \mathbb{I}&\mathbb{O}
\end{pmatrix}.
$$
And the spinor is interpreted as having the first two components representing
\textit{positive energy solutions, particles} and the last two representing
\textit{negative energy solutions, antiparticles}.

Consider now the Dirac equation
$$i\gamma^\mu\pd_\mu\psi(x^\mu) = m\psi(x^\mu)$$
this is clearly a plane wave equation, with solutions in the form
$$\begin{dcases*}
    \psi_1(x^\mu) = u_s(P^i) e^{-iP_\mu x^\mu}\\
    \psi_2(x^\mu) = v_s(P^i) e^{iP_\mu x^\mu}
\end{dcases*}
$$
where $u_s,v_s$ are 4-component spinors that depend only on the $3-$momentum, since
the 0-component is fixed by the dispersion relation
$$E^2 = P_iP^i+m^2$$
which will become important later.

The $\psi_1$ solution is an inbound plane wave with positive energy,
representing a \textbf{particle}, while the other is an outbound
plane wave with negative energy, representing an \textbf{antiparticle}.

The general solution is
$$\psi(x^\mu) = Au_s(P^i)e^{-iP_\mu x^\mu}+ Bv_s(P^i)e^{iP_\mu x^\mu}$$
substituting $\psi_1$ and $\psi_2$ respectively into the equation yields
$$(i\gamma^\mu P_\mu -m)u_s(P^i) = 0 \to (i\slashed{P}-m)u_s(P^i)=0$$
$$(i\gamma^\mu P_\mu +m)v_s(P^i) = 0 \to (i\slashed{P}+m)v_s(P^i)=0$$
Rewriting these equations in the rest frame of the particle, 
$\slashed{P} = \gamma^0P_0 = \gamma^0m$
\begin{gather*}
m(\gamma^0-\mathbb{I})u_s = 0
\iff m\begin{pmatrix}
    \mathbb{O}&\mathbb{O}\\
    \mathbb{O}&-2\mathbb{I}
\end{pmatrix}\begin{pmatrix}
u\\\tilde u
\end{pmatrix} = 0\\
m(\gamma^0+\mathbb{I})v_s = 0
\iff m\begin{pmatrix}
    \mathbb{O}&\mathbb{O}\\
    \mathbb{O}&2\mathbb{I}
\end{pmatrix}\begin{pmatrix}
v\\\tilde v
\end{pmatrix} = 0\\
\end{gather*}
The first equation yields the fact that $\tilde u(0)$ is the $\begin{pmatrix}0\\0
\end{pmatrix}$ vector, and that $u(0)$ is free: so it is chosen to be either
$\begin{pmatrix}1\\0\end{pmatrix}$ or$\begin{pmatrix}0\\1\end{pmatrix}$;
the other equation tells a similar tale, so the $u,v$ are
$$
u_1(0) = \begin{pmatrix}1\\0\\0\\0\end{pmatrix},\;
u_2(0) = \begin{pmatrix}0\\1\\0\\0\end{pmatrix},\;
v_1(0) = \begin{pmatrix}0\\0\\1\\0\end{pmatrix},\;
v_2(0) = \begin{pmatrix}0\\0\\0\\1\end{pmatrix}.
$$
Defining
$$\chi_1 = \begin{pmatrix}1\\0\end{pmatrix},\;
\chi_2 = \begin{pmatrix}0\\1\end{pmatrix}$$
they can be expressed as
$$u_s(0)=\begin{pmatrix}\chi_s\\\mathbf{0}\end{pmatrix},\;
v_s(0)=\begin{pmatrix}\mathbf{0}\\\varphi_s\end{pmatrix}.$$
The normalization is $\bar u u =2m$.

To find the general solution from the rest frame solution, one can perform a
boost to a generic reference frame and get the general solution\footnote{see prof.
Boglione handwritten notes for the explicit calculations.}, but it can also
be found by using the properties of the Clifford Algebra.

Recalling the properties of the Clifford algebra,
$$(\slashed{P}-m)(\slashed{P}+m)=(P_\nu P_\mu\gamma^\mu\gamma^\nu - m^2) = 
(P_\mu P_\nu\frac{1}{2}\underbrace{\left\{\gamma^\mu,\gamma^\nu\right\}}_{
=2g^{\mu\nu}}-m^2)=(P_\mu P^\mu-m^2)=0,$$
so this operator applied to one of the spinors yields again $0$:
\begin{gather*}
(\slashed{P}-m)\underbrace{(\slashed{P}+m)u_s(0)}_{:=u_s(P^i)}=0\to
u_s(P^i)=(\slashed{P}+m)u_s(0)\\
(\slashed{P}+m)\underbrace{(-\slashed{P}+m)v_s(0)}_{:=v_s(P^i)}=0\to
v_s(P^i)=(-\slashed{P}+m)v_s(0).
\end{gather*}
Expressing the parentheses explicitely:
\begin{gather*}
u_s(P^i)=(\slashed{P}+m)u_s(0) = \begin{pmatrix}
    (E+m)\mathbb{I}&\sigma_i P^i\\
    -\sigma_iP^i& -(E-m)\mathbb{I}
\end{pmatrix}\begin{pmatrix}\chi_s\\0\end{pmatrix}=
(E+m)\begin{pmatrix}
\chi_s \\
\frac{-\sigma_iP^i}{E+m}\chi_s
\end{pmatrix}\\
    v_s(P^i) = (E+m)\begin{pmatrix}
    \frac{-\sigma_iP^i}{(E+m)}\\\varphi_s
    \end{pmatrix}. 
\end{gather*}
The orthonormalization relations are
\begin{gather*}
\bar u_s(P^i) u_r(P^i) = 2m\delta_{sr},\\
u^\dagger_s(P^i)u_r(P^i) = 2E\delta_{sr},\\
\bar v_s(P^i)v_r(P^i) = -2m\delta_{sr},\\
v^\dagger_s(P^i)v_r(P^i) = 2E\delta_{sr}.
\end{gather*}
and the spin sums are
\begin{gather*}
    \sum_{s=1}^{2}u_s\bar u_s =\slashed{P}+m,\;
    \sum_{s=1}^{2}v_s\bar v_s =\slashed{P}-m.
\end{gather*}
The energy projectors (useful for something) are
\begin{gather*}
    \left[\Lambda_+(P^i)\right]_{\alpha\beta} =
    \sum_{2=1}^{2} \frac{u^s_\alpha\bar u^s_\beta}{2m}
\end{gather*}
and act like [... anche lambda- vedi foto]

